\section{Related work}\label{sec:related}

\subsection{Weighted ideals, L-shapes, and graph structure}

Zhai's asymptotic theorem \cite{Z} establishes an approximate Wilf
inequality for fixed embedding dimension outside finitely many semigroups.
Our argument instead seeks a uniform, exact result in embedding dimension
four. Within the weighted-ideal framework of Zhai and
Hellus--Rechenauer--Waldi \cite{HRW}, the additional task is a quantitative
surplus sufficient to absorb the genus correction, rather than a new
construction of the exponent ideal.

Aguil\'o-Gost, Garc\'ia-S\'anchez, and Llena \cite{L} study the number
of L-shapes for semigroups of embedding dimension four. Chomicz
\cite{C} gives a geometric procedure for constructing Ap\'ery sets
in that dimension and applies it to explicit families. This is closely
related geometric work, but its relation-deletion and rearrangement
results are not premises of our finite families. We retain the canonical
lexicographic representatives. For any choice of one factorization per
Ap\'ery element, the sum of labels is fixed; rearranging representatives
alone cannot improve that sum. We therefore introduce no optimization
or realizability filter justified only by a change of L-shape.

The clique-tree theorem used in our central-box decomposition is classical
\cite{G}. The application to three-coordinate lower ideals and the resulting
moment bounds are proved explicitly in Section~\ref{sec:analytic-a}.
This use of graph theory is different from Eliahou's graph-theoretic
criterion \(3e\ge m\) \cite{Graph}.

\subsection{Known regions and fixed-multiplicity computation}

Fr\"oberg, Gottlieb, and H\"aggkvist \cite{FGH} prove \(g\le tn\).
Thus \(t\le e-1\) implies
Wilf immediately. Sammartano \cite{Sammartano} proves the criterion
\(2e\ge m\); Eliahou \cite{Graph} improves it to \(3e\ge m\).
Eliahou's Macaulay-theorem argument \cite{Macaulay} settles
\(c\le3m\). These results describe important previously settled regions,
but are not additional assumptions in our branch partition.

The computations of Bruns et al. \cite{B} and Kliem--Stump \cite{KS}
use Kunz cones and their faces. Our small-multiplicity computation instead
enumerates the bounded generator box obtained from the counterexample
reduction in Section~\ref{sec:p6}. This is a different finite assertion,
not a reproduction of the published Kunz-cone computations.

\subsection{Type and presentation bounds}

Chomicz's subsequent paper \cite{CT} bounds the cardinality \(\eta(S)\)
of minimal presentations in terms of the type:
\[
t(S)-11\le\eta(S)\le4t(S)+5.
\]
These bounds concern another aspect of the same Ap\'ery geometry.
They do not give a uniform bound \(t\le3\) for four generators and
therefore do not replace our large-multiplicity branches.
Marashdeh \cite{M} derives type bounds involving Ap\'ery extremal
statistics and reduces Wilf's conjecture to a further inequality.
We do not use those reductions as established substitutes for the
quantitative moment surplus required here.

The cited 2026 works are preprints in the versions listed in the
bibliography. The literature assessment is dated September 17, 2026;
we make no claim that this bibliography certifies publication priority.
